\chapter*{ABSTRACT}
\addcontentsline{toc}{chapter}{ABSTRACT}

Soil salinity, widespread in the Sine-Saloum region, combined with empirical irrigation and fertilization practices, compromises vegetable production. This thesis, entitled ``Optimization of irrigation and fertilization through precision agriculture in saline soils: a comparative study with traditional practice on vegetable crops in Kaolack'', assesses the value of a precision approach in this context.

The experiment, carried out on the Agro TIC experimental field of USSEIN, compares two zones of three 2~m\textsuperscript{2} plots each (turnip, okra and lettuce). The test zone follows an integrated approach: soil study, amendment (biochar, PCA, compost and EGB), then irrigation and fertilization driven by an 8-in-1 sensor, weather data and the FAO-56 Penman-Monteith method. The control zone follows traditional practice, with fixed watering. A progressive web application, AgriSens, designed with the MERISE method and the FIWARE architecture, automates the calculations and the visualization of the data.

The results show that, in the test zone, the soil is looser and more permeable, with no salt efflorescence, along with better growth (okra: 41.3~cm versus 21~cm) and effective production (turnip: 4.8~kg over 2~m\textsuperscript{2}), whereas in the control zone production is limited or even nonexistent. The limitations lie in the estimation of NPK from EC, the absence of ion identification, the low efficacy of biological treatments, and the detection of diseases by simple visual observation.

\vspace{1em}
\noindent \textbf{Keywords:} precision agriculture, saline soils, irrigation, fertilization, evapotranspiration, organic amendment, 8-in-1 sensor, vegetable crops, AgriSens.